% Source: https://artofproblemsolving.com/community/c3h1585825_20th_term
% Stable identity: aops-topic-1585825-q01
The differences of consecutive terms are $5,7,9,11,13,15$: consecutive odd numbers, so the second differences are constant ($2$) and the $n$th term is a quadratic in $n$. Since the differences are $2n+3$ (from term $n$ to term $n+1$),
\[
a_n=13+\sum_{j=1}^{n-1}(2j+3)=13+(n-1)n+3(n-1)=n^2+2n+10 .
\]
Check: $a_1=13$, $a_2=18$, $a_3=25$, $a_7=49+14+10=73$. Hence
\[
a_{20}=400+40+10=\boxed{450}.
\]
(Equivalently, $a_n=(n+1)^2+9$.)
